\tutorial{5}{Restarts, ensembles and particle tracking}

\begin{goal}
With the model of Tutorial~2 you split a run into two and continue the second part from \cmd{solution.rvc}, getting exactly
the continuous run; run a Monte Carlo ensemble whose members rebuild the groundwater model; and write MODFLOW's budget file
to track particles backwards from the pumping well with \cmd{make\_prt.py}. Time: about 30 minutes. The particle tracking
needs FloPy (\cmd{python3 -m pip install flopy}).
\end{goal}

\begin{shell}
$ cp -r $TUT/t5_restarts_ensembles_particles/start ~/work/t5
$ cd ~/work/t5
\end{shell}

\section{Restarts: \texttt{solution.rvc}}
At the end of every run Raven writes the state of the model to \cmd{solution.rvc}: soil and surface stores, channels and
reservoirs, and in a coupled run the head of every MODFLOW cell (and, with a soil zone, its columns). Used as the
\cmd{.rvc} of the next run, it continues the run where it stopped.

\subsection{A continuous run}
The reference: the three years 2010--2012 in one run.
\term{t5_run_full.txt}

\subsection{Part 1: the first two years}
\diff{t5_part1.diff}
\term{t5_run_part1.txt}
\term{t5_rvc_blocks.txt}
After Raven's own blocks (HRU states, basin states) comes \cmd{:GWHeads \V{t5:gwheads}}: the head of every cell, layer by
layer and row by row, with 17 significant digits:
\term{t5_rvc_head.txt}

\subsection{Part 2: the third year, from the hotstart}
Make a folder for the second part with the model's files, and use \cmd{solution.rvc} of part~1 as its \cmd{.rvc}:
\begin{shell}
$ mkdir ../t5_part2
$ cp Creek.* CreekMet.rvt CreekGW.rvt *.geojson *.tif ../t5_part2/
$ cp out_part1/solution.rvc ../t5_part2/Creek.rvc
$ cd ../t5_part2
\end{shell}
(no output). Start on 1 January 2012 and run the rest of the three years:
\diff{t5_part2.diff}
\term{t5_run_part2.txt}
\term{t5_part2_advisory.txt}
The \cmd{:GWInitialization STEADY\_STATE} of the \cmd{.rvg} is skipped: the heads come from the hotstart.

\subsection{Compare with the continuous run}
\term{t5_diff.txt}
\begin{tocheck}
The largest difference is \V{t5:maxdiff}: the restarted run reproduces the continuous one bit for bit.
\end{tocheck}
\begin{trouble}
\begin{itemize}
\item \emph{``running more time steps than the MODFLOW 6 simulation holds''}: you changed \cmd{:Duration} between building
      and running; run again from the start of the folder.
\item A corrupt \cmd{:GWHeads} header or a non-number in a groundwater block stops the run with a message naming it: copy
      \cmd{solution.rvc} again.
\end{itemize}
\end{trouble}
Use restarts to spin a model up once and start every later run from the spun-up state, to split long runs, and for the
unsaturated columns of Tutorial~3, whose contents are saved in \cmd{solution.rvc} as well.

\section{A Monte Carlo ensemble}
Raven's ensemble modes run the model several times; with groundwater, Raven rebuilds the MODFLOW model and restarts
MODFLOW for each member, with its files in the member's output folder. Back in \cmd{\textasciitilde/work/t5}, put the
duration back to three years (undo the part~1 edit) and ask for five Monte Carlo members:
\begin{shell}
$ cd ~/work/t5
\end{shell}
\diff{t5_ens.diff}
The ensemble file \cmd{Creek.rve} says what varies between members: here \cmd{PERC\_COEFF} of the subsoil, which sets how
fast it feeds the aquifer, drawn uniformly between 0.02 and 0.06~1/d. Create it:
\showfile{t5_restarts_ensembles_particles/finished/t5/Creek.rve}
\term{t5_run_ens.txt}
\term{t5_ls_ens.txt}
The run takes about \V{t5:enssec}~seconds (\V{t5:nens} members). Each \cmd{ens\_N} folder holds a member's outputs, with its own
\cmd{mf6/}.
\fig{t5_ensemble.png}{The \V{t5:nens} members: head at OW\_SLOPE and recharge reaching MODFLOW. The members' mean recharge ranges
from \V{t5:ensrmin} to \V{t5:ensrmax}~m\textsuperscript{3}/d and the final head at OW\_SLOPE from \V{t5:enshmin} to
\V{t5:enshmax}~m.}
Remove the two lines of \cmd{:EnsembleMode} and \cmd{:RandomSeed} again before going on (the ensemble file is read only in
ensemble mode). Raven's DDS and EnKF modes rebuild the groundwater model for each member in the same way.

\section{The budget file and particle tracking}
MODFLOW~6's particle tracking (PRT) runs after the coupled run, on the flows it saved. Still in \cmd{\textasciitilde/work/t5}
(with the ensemble lines removed), run one year and ask Raven to save MODFLOW's flows every day:
\diff{t5_prt_rvi.diff}
\diff{t5_prt.diff}
\term{t5_run_prt.txt}
\cmd{out\_prt/mf6/gwf.cbc} now holds every package's flows, the specific discharge and the saturation of every day
(\V{t5:cbcmb}~MB for this year), and \cmd{gwf.porosity} the porosity of every cell (\cmd{POROSITY} of the aquifer classes).
\cmd{make\_prt.py} writes a PRT simulation for these flows --- by default \V{t5:npart} release points around each well of
\cmd{:Wells} --- and runs it with the library; \cmd{--backward} tracks backwards in time, to where the pumped water came
from:
\term{t5_make_prt.txt}
\term{t5_endpoints.txt}
Each line is a particle: when it was released (in days of the tracking), its release point and end point (in metres, in the
frame of the grid file \cmd{gwf.dis.grb}: PRT's own coordinates shifted by its \cmd{XORIGIN} and \cmd{YORIGIN}), the time it was tracked (here the whole year: no particle left the model), the PRT status code of its
last record and, in words, why its track ended. The tutorial's script draws the tracks, in the same frame, with the travel times of \cmd{endpoints.csv}:
\term{t5_plot_tracks.txt}
\fig[0.62]{t5_tracks.png}{Backward tracks from the \V{t5:npart} release points around PW\_1 over one year: the water pumped at the
end of 2010 was \V{t5:distmin} to \V{t5:distmax}~m away at the start of the year.}
\begin{note}[Tracks need small cells]
All the tracks begin and end inside the well's own 250-m cell: PRT interpolates the flow within a cell, and this cell is
itself the sink of the well. They show the method, not a capture zone. A capture-zone or travel-time study needs cells
much smaller than the distances of interest around the well (\cmd{:GridRefinement WELLS}, or a smaller
\cmd{:GridCellSize}).
\end{note}
\begin{trouble}
\begin{itemize}
\item \emph{``No module named flopy''}: \cmd{python3 -m pip install flopy}.
\item \cmd{make\_prt.py} stops with a message when the flows were not saved every day: add \cmd{:HeadSaveFrequency 1} and
      \cmd{:WriteGWBudgetFile} to the \cmd{.rvg}.
\item PRT needs a regular or quadtree grid (DIS or DISV); HRU meshes and imported grids use DISU.
\end{itemize}
\end{trouble}

\section{The finished state}
\cmd{\$TUT/t5\_restarts\_ensembles\_particles/finished} holds the two folders as this tutorial leaves them: \cmd{t5} (with
\cmd{Creek.rve}, the one-year \cmd{Creek.rvi} and the budget-file lines in \cmd{Creek.rvg}) and \cmd{t5\_part2} (with the
hotstart file as \cmd{Creek.rvc}).
